%% ============================================================
%%  SOLUCIÓN — Ejercicio 2.1: Ecuaciones básicas
%%  Clase 1 — Después de diapositiva "Modos matemáticos"
%%  Compilar: F5 (pdflatex, dos veces para referencias)
%% ============================================================

\documentclass[12pt, a4paper]{article}
\usepackage[utf8]{inputenc}
\usepackage[T1]{fontenc}
\usepackage[spanish]{babel}
\usepackage{amsmath}

\title{El modelo econométrico}
\author{Tu Nombre}
\date{\today}

\begin{document}
\maketitle

\section{Modelo de regresión lineal múltiple}

El modelo de regresión lineal múltiple es:
\begin{equation}
  Y_i = \beta_0 + \beta_1 X_{1i} + \beta_2 X_{2i} + \varepsilon_i
  \label{eq:rlm}
\end{equation}

donde $Y_i$ es la variable dependiente, $X_{1i}$ y $X_{2i}$ son los regresores, 
y $\varepsilon_i \sim N(0, \sigma^2)$ es el término de error.

\section{Minimización de la suma de cuadrados}

El estimador OLS minimiza:
\[
  \sum_{i=1}^{N} (Y_i - \hat{Y}_i)^2
    = \sum_{i=1}^{N} (Y_i - \hat{\beta}_0
      - \hat{\beta}_1 X_{1i} - \hat{\beta}_2 X_{2i})^2
\]

Esta es una función cuadrática en los parámetros que se puede minimizar 
usando cálculo. La solución de primer orden nos da las ecuaciones normales 
de OLS.

Como vemos en la ecuación \eqref{eq:rlm}, el modelo es lineal en los 
parámetros. Esta linealidad es crucial porque permite obtener una solución 
cerrada para el estimador OLS.

\end{document}
